\documentclass[aspectratio=169,11pt]{beamer}
\usepackage[T1]{fontenc}
\usepackage{lmodern,booktabs,array}
\definecolor{Ink}{HTML}{17324D}
\definecolor{Teal}{HTML}{007F7B}
\definecolor{Gray}{HTML}{657483}
\definecolor{Rust}{HTML}{B5533C}
\definecolor{Soft}{HTML}{EEF3F6}
\setbeamercolor{alerted text}{fg=Rust}
\setbeamercolor{normal text}{fg=Ink,bg=white}
\setbeamercolor{frametitle}{fg=Ink}
\setbeamercolor{structure}{fg=Teal}
\setbeamercolor{block title}{fg=white,bg=Teal}
\setbeamercolor{block body}{fg=Ink,bg=Soft}
\setbeamertemplate{navigation symbols}{}
\setbeamersize{text margin left=9mm,text margin right=9mm}
\setbeamerfont{frametitle}{size=\large,series=\bfseries}
\setbeamertemplate{footline}{\hspace{9mm}\scriptsize Bootstrapping libraries and their boundaries | advisor update, 24 Oct 2026\hfill\insertframenumber/\inserttotalframenumber\hspace{9mm}\vspace{3mm}}
\newcommand{\source}[1]{\vfill{\scriptsize\color{Gray}#1\par}}
\hypersetup{pdftitle={Where do libraries come from, and where do they stop helping?},pdfauthor={Alex Towell}}
\title{Where do libraries come from, and where do they stop helping?}
\subtitle{Bootstrapped libraries, documentation, evidence selection, and boundaries}
\author{Alex Towell}
\date{Advisor update, 24 October 2026 (revised from 22 October)}
\begin{document}

\begin{frame}[plain]
\titlepage
\vfill
\begin{block}{In one sentence}
Small models handle inputs they cannot read when code holds the work: a 0.6B caller runs an 8B-written MRCR library
to 1M tokens; a trained 4B with three labels met a pre-registered 10-point non-inferiority margin against an 8B with
thirty (75 vs 73 of 141, not under template clustering); bootstrapped libraries break when the question is reworded,
and a 4B question parser repairs them. Recursion is not the active ingredient; acquisition is the open problem.
\end{block}
\end{frame}

\begin{frame}[t]
\frametitle{Libraries can be bootstrapped: an 8B wrote MRCR's library.}
\begin{columns}[T,onlytextwidth]
\begin{column}{.58\textwidth}
\footnotesize
\begin{tabular}{@{}lrr@{}}
\toprule
MRCR held out (8k to 1M tokens) & 2 needles & 4 needles \\
\midrule
Direct answer, 4B & 0.23 & 0.15 \\
RLM, untrained 4B, no library & 0.18 & 0.11 \\
8B-written library, run alone & 1.000 & 1.000 \\
0.6B caller + library & \textbf{1.000} & \textbf{1.000} \\
Untrained 4B + library & 0.95 & 0.88 \\
\bottomrule
\end{tabular}
\medskip

\small 64 rows per column. The 0.6B caller never saw MRCR. The proposer saw the contract and three solved examples; candidates run on 30 teaching tasks; the best scored
30/30 and was reviewed before use (40 lines of parsing, nothing copied from the examples).
\end{column}
\begin{column}{.38\textwidth}
\begin{block}{Who can write it (seeds accepted)}
\small
4B Instruct: 0 of 1\\
4B, thinking off / on: 0 of 1 / 0 of 2\\
8B, thinking off: 0 of 2\\
8B, thinking on: \textbf{2 of 5}\\
14B (4-bit), thinking on: \textbf{2 of 5}\\
4B after SFT + RL: 0 of 4; test-time RL on MRCR: 0 in 1,920 samples
\end{block}
\small \alert{Caveat:} MRCR is solvable by a short parser, so this measures program discovery, not reading. We
found no published note saying so.
\end{column}
\end{columns}
\source{Ledger E21, E21b; runs e21-*. MRCR: openai/mrcr (MIT), official SequenceMatcher grader.}
\end{frame}

\begin{frame}[t]
\frametitle{Discovery on 37 generated families: thinking, scale, and training.}
\begin{columns}[T,onlytextwidth]
\begin{column}{.55\textwidth}
\footnotesize
\begin{tabular}{@{}lrr@{}}
\toprule
Proposer & Solved / 37 & Tokens \\
\midrule
4B Instruct (2 seeds) & 5, 4 & 4k \\
8B, no thinking & 4 & 4k \\
4B, thinking & 14 & 55k \\
8B, thinking (2 seeds) & 20, 25 & 48k \\
14B, thinking (2 seeds) & 21, 26 & 30k \\
4B + SFT (2 seeds) & 9, 7 & 3k \\
4B + SFT + RL (2 seeds) & 16, 14 & 2k \\
4B, round four (2 seeds) & \textbf{22, 20} & 2k \\
\quad avg. of 4 RL seeds & 15.8 & \\
4B, 3x helper data, avg. & 16.8 & 2k \\
\bottomrule
\end{tabular}
\medskip

\small Tokens: proposer output per family. Families vary record format, question and answer format; held-out
families use a new format or question.
\end{column}
\begin{column}{.41\textwidth}
\small
\textbf{Training the proposer:} SFT on the 8B's verified code, then 40 steps of RL with the verified teaching
score as reward: 9 to 30 of 74 attempts (p $<$ 0.001). \alert{Seed-sensitive:} round four's 42 (8B 45) is the
best of four RL seeds (22 to 42). MRCR still fails.\par\medskip
\textbf{Helper gap:} 3x the helper teacher data matches round four on solves, reaches the 8B on helper families
(mean 0.46), and varies far less by seed.
\end{column}
\end{columns}
\source{Ledger E26 to E26g, E36 to E46; runs e26-disc-*, e4*. Generator: family\_gen.py (answers exact).}
\end{frame}

\begin{frame}[t]
\frametitle{Bootstrapping is cheap where answers are exact (pre-registered).}
\begin{columns}[T,onlytextwidth]
\begin{column}{.56\textwidth}
\footnotesize
\begin{tabular}{@{}lrrr@{}}
\toprule
 & Correct & False acc. & Tokens \\
\midrule
\multicolumn{4}{@{}l}{\emph{E39: 30 new families, 2 seeds}} \\
4B, best of 8, 3 labels & 12, 14 & 0, 1 & 3.7k \\
8B thinking, 30 labels & 13, 16 & 0, 0 & 52.4k \\
\multicolumn{4}{@{}l}{\emph{E58: 141 untouched families (powered)}} \\
4B, best of 8, 3 labels & \textbf{75} & \alert{10} & 21k \\
8B thinking, 30 labels & 73 & 0 & 64k \\
\bottomrule
\end{tabular}
\medskip

\scriptsize Correct = 0.9+ on 30 fresh tasks; tokens per family (E39: output only). E58 met the registered 10-point
non-inferiority margin with families as the unit (95\% CI $-0.078$ to $+0.106$), not with templates as the unit; the
weakest checkpoint did not (67 vs 73).
\end{column}
\begin{column}{.40\textwidth}
\begin{block}{Why three labels suffice}
\small
Candidates that run cleanly and reproduce the prompt's examples are correct: precision 0.81, 0.88, 0.98 with 1, 2,
3 labels (127 of 129 at three, all generalizing). Agreement among candidates is no substitute: an RL-trained
proposer repeats its wrong programs.
\end{block}
\small \alert{Limits:} noisy families need tens of labels; E58's false acceptances: all 6 semantic totals, 3 of 7
day spans.
\end{column}
\end{columns}
\source{Ledger E37 to E39, E58; notebook assets-e39- and assets-e58-preregistration.md.}
\end{frame}

\begin{frame}[t]
\frametitle{Libraries generalize in length, not in wording; the harness can fix that.}
\begin{columns}[T,onlytextwidth]
\begin{column}{.52\textwidth}
\footnotesize
\begin{tabular}{@{}lrrr@{}}
\toprule
Unseen shift (of 55) & As is & 4B & 0.6B \\
\midrule
Inputs 5$\times$ longer & 55 & & \\
Question reworded & 3 & \textbf{54} & 22 \\
Prefix, field order & 11 & \textbf{48} & 26 \\
\bottomrule
\end{tabular}
\medskip

\small The parser wrapper: a 4B reads only the question, returns the family's parameters (96\% exact on unseen
rewordings), and the question is re-rendered in the original template; the frozen library runs unchanged.
Remaining failures are structured formats whose records changed.
\end{column}
\begin{column}{.44\textwidth}
\begin{block}{Who can supply the robust interface?}
\small
4B proposer, told how, writing from scratch: 1 of 11 and 0 of 18 accepted libraries adopt it.\\
4B editing a working library: working edits for 3 of 12 families (2 robust).\\
Hand-written wrapper + 4B parser: 54 of 55.
\end{block}
\small Executors generalize; a 4B parses well; a small proposer rarely writes the interface. \alert{Acquisition is
the open problem.}
\end{column}
\end{columns}
\source{Ledger E59 to E61, E68, E69; stress\_gen.py, e68\_canon.py.}
\end{frame}

\begin{frame}[t]
\frametitle{One line of documentation decides whether an untrained caller uses a library.}
\begin{columns}[T,onlytextwidth]
\begin{column}{.5\textwidth}
\footnotesize
\begin{tabular}{@{}lr@{}}
\toprule
MRCR, untrained 4B caller, docstring & Score \\
\midrule
Signature only & 0.18 \\
Description + signature (as the 8B wrote it) & 0.59 \\
Signature + usage line & 0.95 \\
Description + signature + usage line & 0.95 \\
\midrule
Hand-written library, usage line removed & 0.22 \\
Hand-written library, as written & 0.95 \\
\bottomrule
\end{tabular}
\end{column}
\begin{column}{.46\textwidth}
\small
Usage line: \texttt{Typical use: FINAL\_TEXT(answer\_mrcr(request)).}\par\medskip
Without it the untrained 4B hands the whole input to a sub-model (overflowing at 128k and above) or retypes the
routine's output by hand.\par\medskip
\textbf{The distilled 0.6B scores 1.000 without the line; the untrained 0.6B fails either way (0.21, 0.13).}
Robustness comes from distillation, not size. The bootstrapper now inserts the line.
\end{column}
\end{columns}
\source{Ledger E25c, E25d. One family, one library, one seed per arm.}
\end{frame}

\begin{frame}[t]
\frametitle{MuSiQue in long haystacks: capacity belongs in choosing evidence.}
\centering\footnotesize
\begin{tabular}{@{}lrrrr@{}}
\toprule
Answer F1 (90 to 180 held-out questions per row) & 8k & 32k & 128k & tokens \\
\midrule
Direct answer, 4B & 0.52 & 0.25 & overflow & \\
RLM without a library, depth 0 / depth 1 & 0.49 / 0.51 & 0.25 / 0.25 & overflow & \\
Library that hands each hop to a child RLM & 0.29 & 0.29 & 0.34 & 3 to 4 times \\
0.6B caller + filter, BM25, read & 0.52 & 0.31 & 0.46 & 57k \\
0.6B caller + two-round retrieval & \textbf{0.53} & \textbf{0.39} & \textbf{0.50} & 10k \\
\bottomrule
\end{tabular}
\par\medskip\raggedright
\begin{columns}[T,onlytextwidth]
\begin{column}{.48\textwidth}
\small \textbf{What helps:} a second, targeted retrieval (BM25 for the question, the helper names the missing
fact, BM25 again, one read): +0.058 over BM25 alone, +0.045 over filter+BM25 at a sixth of the tokens.
\end{column}
\begin{column}{.48\textwidth}
\small \textbf{What does not:} recursion; an 8B reader (+0.01; gold paragraphs reach 0.63 to 0.67 with either
reader); more sampling (majority of 3 at 88k tokens: 0.22 vs 0.35 for the library at 32k).
\end{column}
\end{columns}
\source{Ledger E22c to E22f, E23, E25e, E28. MuSiQue (CC BY 4.0) dev questions in haystacks.}
\end{frame}

\begin{frame}[t]
\frametitle{External validity, and what we withdrew.}
\begin{columns}[T,onlytextwidth]
\begin{column}{.48\textwidth}
\begin{block}{LongBench-v2 (all 503 questions, chance 0.25)}
\footnotesize
Direct (truncated): 0.356, 42k tokens\\
\textbf{Retrieve-then-read: 0.427, 10k tokens}\\
via the 0.6B caller: 0.423 (+0.07, se 0.02)\\[1mm]
Selecting evidence beats summarizing (map-reduce: no gain).
\end{block}
\footnotesize \textbf{Boundary:} no proposer bootstraps OOLONG, our own family (8B 0.20, 14B 0.00; shipped
library 0.60): 30 examples span 11 question templates.\par\smallskip
\textbf{Withdrawn:} distractor-limited reader; length routing (pre-registered); ``no help on
LongBench-v2'' (it was map-reduce).
\end{column}
\begin{column}{.48\textwidth}
\footnotesize
\textbf{Method lessons}\par\smallskip
\textbullet\ Verify under the semantics you run under: the kernel's \texttt{ask} raises on a cut-off reply,
scripted harnesses returned partial text (one routine crashed in 46 of 90 caller runs).\par\smallskip
\textbullet\ Verification must span deployment lengths: 8k-only teaching selected a pass-through routine that
scores 0 at 128k.\par\smallskip
\textbullet\ Routine evolution (GEPA-style): prompt rewrites by the 8B hurt; numbers helped (+0.048, ns);
random search found nothing.
\end{column}
\end{columns}
\source{Ledger E32 to E32e, E31, E22b, E24, E24c, E35. LongBench-v2: zai-org (Apache-2.0).}
\end{frame}

\begin{frame}[t]
\frametitle{What this changes, and what to test next.}
\small
\textbf{The thesis, revised.} For small models on large inputs, capability comes from code with programmatic access
and tested routines, plus narrow neural decisions (parse a question, classify an item). Recursion is not the active
ingredient; we keep the term RLM and frame it as a harness. The mechanisms are prior art (CodeAct, RLM,
skill libraries, FlashExtract, Evaporate); what is new is measured on small models: \alert{acquiring} robust routines
cheaply, and where that breaks (wording, lengths beyond teaching, false acceptance).\par\medskip
\textbf{Next}
\begin{enumerate}\setlength\itemsep{1mm}
\item The deciding experiment: interface (tool loop vs REPL) $\times$ source of procedures (primitives,
human-written, acquired), plus a host pipeline with no controller. Is the thesis about the harness or acquisition?
\item False acceptance without labels: metamorphic checks (shift times, keep dates) and helper consistency.
\item Who writes the interface: an 8B proposer (running), or the harness supplies it.
\item Acceptance must span deployment lengths and wordings (E52b, E57, E59).
\end{enumerate}
\end{frame}

\begin{frame}[t]
\frametitle{Appendix: protocol and provenance.}
\footnotesize
Models: Qwen3-4B-Instruct-2507 (helper, caller), Qwen3-0.6B + LoRA (distilled caller), Qwen3-4B hybrid and
Qwen3-8B-AWQ (proposers), all served with vLLM 0.27.1 on two A100 MIG 7g.40gb slices.\par\smallskip
Data: OpenAI MRCR (MIT; within-bin split, 30 teaching rows, 64 held out per needle count); MuSiQue v1.0 (CC BY
4.0; teaching from train, held out from dev; two samples of 90); LongBench-v2 (Apache-2.0, revision 2b48e494);
generated families (family\_gen.py, exact answers).\par\smallskip
Statistics: paired differences with standard errors over all questions; exact McNemar for family counts;
the routing test was pre-registered.\par\smallskip
Every number above has a ledger row in docs/RESULTS.md and a run directory in the run store; the lab notebook
(docs/notebook/2026-10-19-...) records failures and corrections in time order.
\end{frame}

\end{document}
